% v13 research notes: reusable local classification, not an unresolved C5 claim.
\subsection{二つの負の reverse divisors が四以下の完全な局所 catalog}
\label{nr13:bounded-reverse}

二高分母標準形の負の高低 entries は $-s,-y$ であり，
それぞれの reverse numerators の絶対値は $S,r+1$ である．
これらがともに四以下という局所条件を課すと，
分母の上限を置かずに全ての整数模型を有限表と三つの族に分けられる．
この catalog は，今後別の support に現れる共通 $\lambda=0$ punctures を調べる際に再利用できる．
五角形成分自体は v13 で排除されており，その残存候補を意味しない．

\paragraph{仮定と完全性．}
正の整数二高分母標準形で $a,b$ は奇数とし，
\[
 2\leq S\leq4,\qquad 1\leq r\leq3
\]
を仮定する．$S\geq2$ は負の低分母 entry の真約数性から従う．
このとき，次の表の二十二模型と三つの無限族だけが可能である．
表は canonical data の分類であって，標識を忘れた幾何的同型類を数えるものではない．

\begin{center}
\small
\begin{tabular}{rrrrrrrrr}
$r$&$S$&$T$&$s$&$z$&$a$&$b$&$D_A$&$D_B$\\\hline
1&2&5&1&2&3&11&6&$-10$\\
1&2&5&2&6&5&31&18&$-27$\\
1&3&8&2&5&7&41&21&$-36$\\
1&3&9&2&2&7&19&7&$-15$\\
1&3&9&4&6&13&55&23&$-41$\\
1&3&9&6&26&19&235&127&$-211$\\
1&3&11&8&2&25&23&$-9$&9\\
1&3&12&4&1&13&13&$-3$&$-2$\\
1&4&9&1&4&5&37&21&$-36$\\
1&4&10&1&2&5&21&11&$-21$\\
1&4&12&1&1&5&13&6&$-15$\\
1&4&12&3&5&13&61&26&$-49$\\
1&4&12&4&13&17&157&78&$-138$\\
1&4&13&9&28&37&365&177&$-312$\\
1&4&13&10&62&41&807&428&$-747$\\
1&4&15&6&2&25&31&$-4$&$-1$\\
1&4&15&25&4&101&61&$-63$&102\\
1&4&16&3&1&13&17&0&$-9$\\
1&4&16&13&2&53&33&$-31$&45\\
1&4&18&7&1&29&19&$-15$&15\\
2&3&5&2&6&7&31&21&$-25$\\
2&4&7&3&12&13&85&58&$-73$
\end{tabular}
\end{center}

無限族は $t,k\geq1$ を整数として
\[
\begin{array}{c|c|c|c}
 (r,S,T,s,z)&(a,b)&(D_A,D_B)&\text{parameter}\\\hline
 (1,2,6,2t-1,t)&(4t-1,6t+1)&(3,-7)&t\geq1\\
 (2,3,6,2k,1)&(6k+1,7)&(6-6k,7k-9)&k\geq1\\
 (3,4,6,3k,2)&(12k+1,13)&(12-12k,13k-14)&k\geq1
\end{array}
\]
となる．第一族では二つの分母がともに増大する．
後の二族は固定分母七・十三の共鳴であり，公開稿の
\texttt{sp13:no-resonant-continuation} により，固定 ray をもう一つの二高分母 core に
接続して対向 $D$ を相殺することはできない．

\paragraph{導出：小さい reverse data から有限個の factorization へ．}
$h=S-r$，$k=T-r-1$ と置く．$h=0$ は canonical 補助変数 $\alpha$ の非零性に反する．
$h<0$ の可能性は $(r,S)=(3,2)$ だけである．
負の枝の二分により $T=S=2$ となり，整除性から $s=3u,z=2v$，$u,v\geq1$ と書ける．
このとき $\alpha=-u,\beta=-v$ であり，Bézout 式は
\[
 (9u+1)(-v)-(8v+1)(-u)=1,
 \quad\text{すなわち}\quad u(1-v)-v=1
\]
となるが，左辺は負である．従って六組
\[
 (r,S)=(1,2),(1,3),(1,4),(2,3),(2,4),(3,4)
\]
だけを調べればよい．全て $h,k>0$ である．

\[
 R=r(r+1),\quad U=(r+1)(S-r),\quad V=r(T-r-1),
 \quad E=r^2T-(r+1)^2S+R
\]
と定義すると，高高 pair product は
\begin{equation}\label{nr13:bilinear}
 Esz+Vz-Us=R
\end{equation}
と同値である．$s,z\geq1$ より
\begin{equation}\label{nr13:T-bound}
 s(E-U)\leq R-V
\end{equation}
である．一方，正の枝の不等式 $rT>(r+1)S$ を保持する．
これらから $T$ は次の二十六候補に限られる：
\[
\begin{array}{c|c}
 (r,S)&T\\\hline
 (1,2)&5,6,7\\
 (1,3)&7,8,9,10,11,12,13\\
 (1,4)&9,10,11,12,13,14,15,16,17,18,19\\
 (2,3)&5,6\\
 (2,4)&7,8\\
 (3,4)&6
\end{array}
\]
実際，$r=1$ で $T\geq6S-4$ なら $E-U\geq0$，$R-V=4-T<0$ となり矛盾する．
$(r,S)=(2,3)$ で $T\geq7$，$(2,4)$ で $T\geq9$ も同様である．
$(3,4)$ で $T\geq7$ なら
$s(E-U)\geq9T-56>24-3T=R-V$ となって矛盾する．

$E\ne0$ のとき，式~\eqref{nr13:bilinear} を
\begin{equation}\label{nr13:factorization}
 (Es+V)(Ez-U)=ER-UV
\end{equation}
と書く．$C=ER-UV\ne0$ なら $d=Es+V>0$ は $|C|$ の正約数であり，
\[
 s=\frac{d-V}{E},\qquad z=\frac{Us+R}{d}
\]
を調べればよい．正整数性，$r\mid sS$，$r+1\mid zT$，$a=sS+1,b=zT+1$ の奇数性を課すと，
表示した二十二模型がちょうど残る．候補生成と棄却理由は
\texttt{bounded\_reverse\_catalog\_verification\_v13.json} に全て記録している．

$E=0$ の三候補は $r=1,T=4S-2$ であり，式は
$(S-1)(2z-s)=1$ となる．従って $S=2,s=2z-1$ だけが残り，第一族を得る．
$E\ne0,C=0$ の二候補は $(r,S,T)=(2,3,6),(3,4,6)$ であり，
式~\eqref{nr13:factorization} と $Es+V>0$ から，それぞれ $z=1,z=2$ となる．
残る整除性が $s=2k,s=3k$ を与え，後の二族を得る．
全ての模型は元の pair products と row sums を満たし，
標準形の明示 weights により正である．

\paragraph{この分類の射程．}
有限表の最大分母八〇七は完全な factorization の結果として現れた数であり，
計算時の探索上限ではない．ただし reverse divisors の上限四は仮定である．
他の問題へ適用するときは，その上限が本当に node 条件等から従うことを先に示す必要がある．
また局所 charts の表は，整数 marked clutching，node 配分，actual areas の大域的実現を保証しない．
